% ===========================================================================
%  Hans Halvorson --- curriculum vitae
%
%  BUILD.  site.hs (`buildCv`) runs latexmk on this file and serves the result
%  as /cv.pdf.  The two \input fragments below, cv-talks.tex and
%  cv-outreach.tex, are GENERATED from data/talks-master.yaml on every build.
%  Never edit them by hand --- edit the YAML.
%
%  PUBLICATIONS.  The list is driven by hh.bib.  Entries route themselves into
%  the right subsection; only four kinds of entry need a `keywords` tag:
%
%      keywords = {book-authored}  -> Books
%      keywords = {book-edited}    -> Edited volumes
%      keywords = {review}         -> Book reviews
%      keywords = {thesis}         -> suppressed (PhD appears under Education)
%      (untagged)                  -> Articles if @article, else Chapters
%
%  So a new article or chapter needs no tag at all.  Numbering runs backwards
%  within each subsection (newest = highest), which requires knowing each
%  subsection's size in advance; the \AtDataInput hook below counts them.
%
%  SELF-AUTHORSHIP.  The sourcemap strips "Halvorson" from author and editor
%  lists, so sole-authored work prints bare and joint work prints "(with ...)".
% ===========================================================================

\documentclass[11pt]{article}
\usepackage[top=1in,bottom=1in,left=1in,right=1in]{geometry}

% --- Bibliography ----------------------------------------------------------
\usepackage[style=phys,backend=biber,doi=true,sorting=ydnt]{biblatex}
\addbibresource{/Users/hhalvors/hhalvors.github.io/hh.bib}

% Per-subsection entry counts, tallied as biber's data is read in.  Order of
% the tests matters: keyword tests first, then fall back to entry type.
\newcounter{npubbook}
\newcounter{npubedited}
\newcounter{npubarticle}
\newcounter{npubchapter}
\newcounter{npubreview}
\AtDataInput{%
  \ifkeyword{thesis}
    {}%
    {\ifkeyword{review}
       {\stepcounter{npubreview}}%
       {\ifkeyword{book-authored}
          {\stepcounter{npubbook}}%
          {\ifkeyword{book-edited}
             {\stepcounter{npubedited}}%
             {\ifentrytype{article}
                {\stepcounter{npubarticle}}%
                {\stepcounter{npubchapter}}}}}}}

% Reverse numbering.  \pubsfrom{<counter>} primes the countdown; each \item
% decrements it, so the first (newest) entry printed carries the total.
\newcounter{bibnum}
\newcommand*{\bibnumlabel}{\number\value{bibnum}.}
\newcommand*{\pubsfrom}[1]{\setcounter{bibnum}{\numexpr\value{#1}+1\relax}}

\defbibenvironment{bibliography}
  {\list
     {\bibnumlabel}
     {\settowidth{\labelwidth}{99.}%       box exactly fits the widest 2-digit number,
      \setlength{\labelsep}{0.5em}%        so double-digit numbers sit flush at the margin
      \setlength{\leftmargin}{\labelwidth}%
      \addtolength{\leftmargin}{\labelsep}%
      \setlength{\itemsep}{\bibitemsep}%
      \setlength{\parsep}{\bibparsep}%
      \renewcommand*{\makelabel}[1]{\hss##1}}}% \hss = right-align (## required here)
  {\endlist}
  {\addtocounter{bibnum}{-1}\item}

% biblatex-phys drops the DOI into the volume/pages slot whenever an entry has
% no page range --- i.e. on every forthcoming article --- and then prints it
% again at the end of the entry.  Suppress the first, duplicated occurrence.
\renewbibmacro*{note+pages}{%
  \printfield{note}%
  \setunit{\bibpagespunct}%
  \printfield{pages}}

\DeclareSourcemap{
  \maps[datatype=bibtex]{
    \map{
      \step[fieldsource=author,
            match=\regexp{(\w+\s+Halvorson|Halvorson,\s*\w+|H\.\s+Halvorson|Halvorson,\s*H\.)},
            replace={HH}]
      \step[fieldsource=editor,
            match=\regexp{(\w+\s+Halvorson|Halvorson,\s*\w+|H\.\s+Halvorson|Halvorson,\s*H\.)},
            replace={HH}]
      % Then strip "HH": sole-authored works become author-less; co-authored
      % works keep only the co-authors (printed as "(with ...)" by the macro below).
      \step[fieldsource=author, match=\regexp{\AHH\s+and\s+}, replace={}]% HH first
      \step[fieldsource=author, match=\regexp{\s+and\s+HH\b}, replace={}]% HH middle or last
      \step[fieldsource=author, match=\regexp{\AHH\Z},        replace={}]% HH sole author
      \step[fieldsource=editor, match=\regexp{\AHH\s+and\s+}, replace={}]
      \step[fieldsource=editor, match=\regexp{\s+and\s+HH\b}, replace={}]
      \step[fieldsource=editor, match=\regexp{\AHH\Z},        replace={}]
    }
  }
}

% --- Author display ---------------------------------------------------------
% HH is removed from the author list (sourcemap above), so:
%   * sole-authored works have an empty author -> print nothing;
%   * co-authored works keep only the co-authors -> print "(with ...)".
\renewbibmacro*{author}{%
  \ifnameundef{author}
    {}%
    {\mkbibparens{with\addspace\printnames[given-family]{author}}}}
% Same treatment for the byline of a volume HH edited: co-editors print as
% "(with ...), eds.", a volume he edited alone prints nothing at all.  This is
% the byline macro only --- `byeditor+others', which sets the "edited by ..."
% of a volume containing one of HH's chapters, is deliberately left alone.
\renewbibmacro*{editor+others}{%
  \ifnameundef{editor}
    {}%
    {\mkbibparens{with\addspace\printnames[given-family]{editor}}}%
  \clearname{editor}}% or phys reprints them as "edited by ..." further along

% --- Layout ----------------------------------------------------------------
\usepackage{xurl}
\usepackage[none]{hyphenat}
% Hyphenation is off, so TeX has only interword glue to work with.  Without
% extra stretch, long proper nouns and grant titles push into the margin.
\setlength{\emergencystretch}{3em}
\usepackage{microtype}
\usepackage{xcolor}
\definecolor{darkblue}{rgb}{0,0,0.5}

\usepackage{enumitem}
% One list style for the whole document: small mid-dot, tight indent.  Nested
% lists get an en dash.  Everything below uses plain `itemize`.
\setlist[itemize,1]{label=\textperiodcentered, leftmargin=1.1em,
                    itemsep=0.15\baselineskip, topsep=3pt, parsep=0pt}
\setlist[itemize,2]{label=\textendash, leftmargin=1.2em,
                    itemsep=0.1\baselineskip, topsep=2pt, parsep=0pt}

\usepackage{marginnote}
\reversemarginpar
\newcommand{\yearlabel}[1]{\marginnote{\raggedleft\textcolor{gray}{\small\bfseries#1}}}

\usepackage{titlesec}
\titleformat{\subsection}[block]{\normalfont\large\scshape}{}{0em}{}[\titlerule]
\titleformat{\subsubsection}[block]{\normalfont\bfseries}{}{0em}{}
\titlespacing*{\subsubsection}{0pt}{1.1em}{0.35em}

\usepackage{fancyhdr}
\pagestyle{fancy}
\renewcommand\headrulewidth{0pt}
\renewcommand\footrulewidth{0pt}
\lhead{}\rhead{}
\usepackage{lastpage}
\cfoot{\small p \thepage\ of \pageref*{LastPage}}

\setlength{\parindent}{0em}
\newcommand{\journal}[1]{\textit{#1}}
% Dated entries: year in the left column, item in the right.
\newenvironment{dated}[1][ll]%
  {\begin{tabular}{@{}#1}}%
  {\end{tabular}}
% Sub-heading used inside Publications, Teaching and Service.  Same shape as
% \subsubsection* but usable where a numbered sectional unit would be wrong.
\newcommand{\rubric}[1]{\subsubsection*{#1}}

\usepackage{hyperref}
\hypersetup{colorlinks=true,linkcolor=cyan,urlcolor=darkblue,
            pdfborderstyle={/S/U/W 1},
            pdftitle={Hans Halvorson --- Curriculum Vitae},
            pdfauthor={Hans Halvorson}}

\begin{document}
\thispagestyle{empty}

\vspace*{-4em}
\begin{center}
  {\Large\sc Curriculum Vit\ae} \\ \smallskip
  {\Large\sc Hans Halvorson} \\ \smallskip
  {\normalsize Stuart Professor of Philosophy, Princeton University}
\end{center}

\vspace{1em}
\begin{tabular}{@{}l}
  Department of Philosophy, Princeton University \\
  Princeton, NJ 08544 \\
  \href{https://hanshalvorson.com}{hanshalvorson.com}
\end{tabular}


\subsection*{Appointments}

\begin{dated}
  2016--   & Stuart Professor of Philosophy, Princeton University \\
  2012--   & associated faculty, Department of Mathematics, Princeton University \\
  2010--   & professor of philosophy, Princeton University \\
  2005--10 & associate professor, Princeton University \\
  2001--05 & assistant professor, Princeton University
\end{dated}

\subsection*{Visiting appointments and affiliations}

\begin{dated}
  2021--   & research professor, University of Copenhagen \\
  2008--09 & visiting researcher, Mathematical Institute, Utrecht \\
  2005--06 & fellow, Perimeter Institute for Theoretical Physics \\
  2002--   & fellow, Pittsburgh Center for the Philosophy of Science
\end{dated}

\subsection*{Education}

\begin{dated}
  2001 & PhD philosophy, University of Pittsburgh \\
       & \hspace{1em}\emph{Locality, localization, and the particle concept:} \\
       & \hspace{1em}\emph{topics in the foundations of quantum theory} \\
  1999 & visiting student researcher, Oxford University \\
  1998 & MA mathematics, University of Pittsburgh \\
  1997 & MA philosophy, University of Pittsburgh \\
  1995 & BA philosophy, with honors, Calvin College
\end{dated}


\subsection*{Honors and prizes}

% Reverse chronological.  Fellowships that carried research money are listed
% under Research funding instead, so nothing is counted twice.
\begin{itemize}
\item Behrman Fellow, Princeton Council of the Humanities, 2005--08
\item \href{https://www3.nd.edu/~cushpriz/}{James Cushing Memorial Prize in
  the History and Philosophy of Physics}, 2004
\item ``Ten best articles in philosophy for the year 2002,''
  \textit{The Philosopher's Annual}
\item ``Best article by a recent PhD for the year 2001,'' Philosophy of
  Science Association
\item ``Ten best articles in philosophy for the year 2001,''
  \textit{The Philosopher's Annual}
\item Alan Ross Anderson Fellowship for Philosophical Logic, University of
  Pittsburgh, 1997
\item Pew Younger Scholars Graduate Fellow, 1995--98
\end{itemize}

\subsection*{Research funding}

\begin{itemize}
\item Princeton Council of the Humanities grant, ``The Modern Breakthrough in
  Scandinavia: Philosophy, Science, Art'' (co-PI Bridget Alsdorf), 2024--26.
  \$75,000
\item John Templeton Foundation research grant, ``Space from entanglement,''
  2024--26. \$225,000
\item Templeton World Charity Foundation research grant, ``Experimental tests
  of quantum reality'' (co-PI Andrew Briggs), 2013--16. \$2,600,000
\item National Science Foundation research grant, ``Identical particles and
  statistics in superselection theory'' (co-PI David J.\ Baker), 2011--12.
  \$100,000
\item Mellon New Directions in the Humanities Fellowship, 2008--09. \$250,000
\item John Templeton Foundation research grant, ``Understanding the quantum
  world through mathematical innovation,'' 2007
\item Princeton Tuck Grant for International Research Collaboration, 2018--19
  (Denmark); 2010--11 (Netherlands)
\item Princeton Center for the Study of Religion grant for developing the
  course ``Religion and Scientific Objectivity,'' 2017
\item Princeton 250th Anniversary Fund for Innovation in Undergraduate
  Education, grant for developing the logic curriculum, 2004
\end{itemize}


\subsection*{Publications}

\nocite{*}
\renewcommand{\bibsetup}{\setlength{\bibitemsep}{0.5em}\raggedright}

\subsubsection*{Books}
\pubsfrom{npubbook}
\printbibliography[heading=none,keyword=book-authored]

\subsubsection*{Edited volumes}
\pubsfrom{npubedited}
\printbibliography[heading=none,keyword=book-edited]

\subsubsection*{Articles}
\pubsfrom{npubarticle}
\printbibliography[heading=none,type=article,notkeyword=review]

\subsubsection*{Chapters in books}
\pubsfrom{npubchapter}
\printbibliography[heading=none,type=incollection]

\subsubsection*{Book reviews}
\pubsfrom{npubreview}
\printbibliography[heading=none,keyword=review]

% Keep in step with drafts.md, which drives /drafts.html.
\subsubsection*{Unpublished papers and notes}

\begin{itemize}

\item ``Invariance and ontology in relativistic physics''
  \href{https://philpapers.org/rec/HALTIN-3}{philpapers:HALTIN-3}

\item ``Momentum and context''
  \href{https://philpapers.org/rec/HALMAC-2}{philpapers:HALMAC-2}

\item ``Quantifier variance without collapse''
  \href{https://philpapers.org/rec/HALQVW}{philpapers:HALQVW}

\item ``Against eliminating sorts''
  \href{https://philpapers.org/rec/HALAES-4}{philpapers:HALAES-4}

\item ``Mind independence and invariance''

\item ``Does quantum theory kill time?''
  \href{https://philpapers.org/rec/HALDQT}{philpapers:HALDQT}

\item (with S. Wolters) ``Independence conditions for nets of local algebras
  as sheaf conditions''
  \href{http://arxiv.org/abs/1309.5639}{arxiv:1309.5639}

\item (with N. Swanson) ``Comment on the structure of physics''
  \href{http://philsci-archive.pitt.edu/9314}{philsci-archive:9314}

\end{itemize}


\subsection*{Writing for wider audiences}

\begin{itemize}

\item ``Niels Bohr forvansket''
  \href{https://www.weekendavisen.dk/2020-23/ideer/debat}{Weekendavisen,
    Jun 2020}

\item ``Can science speak about the past''
  \href{https://biologos.org/blogs/guest/can-science-speak-about-the-past}{Biologos,
    Dec 2017}

\item ``Fine-tuning does not imply a fine tuner''
  \href{http://cosmos.nautil.us/short/119/fine-tuning-does-not-imply-a-fine-tuner}{Nautilus,
    Jan 2017}

\item ``Matter'' \href{https://www.edge.org/response-detail/27030}{Edge Question 2017}

\item ``Einstein was wrong'' \href{https://edge.org/response-detail/26748}{Edge Question 2016}

\item ``What does quantum mechanics suggest about our perceptions of
  reality?''
  \href{https://www.bigquestionsonline.com/content/what-does-quantum-mechanics-suggest-about-our-perceptions-reality}{Big
    Questions Online 2015}

\item ``When is belief in miracles rational?''
  \href{http://www.slate.com/bigideas/are-miracles-possible/essays-and-opinions/hans-halvorson-opinion}{Slate
    2015}

\item ``Meta-thinking'' \href{http://edge.org/response-detail/26188}{Edge Question 2015}

\item (with Adam Neder) ``Quantum envy''
  \href{http://www.lutheranforum.org/categories/archive/winter-2007/LF2007-4_55-57_Halvorson-Neder-Quantum_Envy.pdf}{Lutheran
    Forum 2007}

\item ``Comments on Clouser's claims for theistic science''
  \href{http://www.asa3.org/ASA/PSCF/2006/PSCF3-06Halvorson.pdf}{PSCF 2006}

\end{itemize}


\subsection*{Invited lectures and conference presentations}

\begin{itemize}
\input{cv-talks.tex}
\end{itemize}

\subsection*{Public lectures and outreach}

\begin{itemize}
\input{cv-outreach.tex}
\end{itemize}


\subsection*{Teaching and advising}

\rubric{Doctoral and postdoctoral advising}

I have worked with and written letters for the PhD students and postdocs
listed below --- ten of them as primary dissertation advisor, marked with an
asterisk. Where a student's PhD is from another university, that university is
given in parentheses. Unless otherwise indicated, the placement listed is a
tenure-track position.

\begin{itemize}
\item Jake Khawaja*
\item Brendan Kolb*
\item Alexander Meehan*: postdoc Yale, University of Wisconsin at Madison
\item David Schroeren*: postdoc Geneva
\item Daniel Berntson: postdoc Rutgers, postdoc Uppsala
\item Robbie Hirsch*: postdoc Princeton
\item Thomas Barrett*: postdoc NYU, UC Santa Barbara
\item Dimitris Tsementzis*: postdoc Rutgers, Goldman Sachs machine learning quants
\item Laurenz Hudetz (Salzburg): London School of Economics
\item Neil Dewar (Oxford): postdoc Munich, Cambridge University
\item Michaela McSweeney: Boston University
\item Benjamin Feintzeig (UC Irvine): University of Washington
\item Noel Swanson*: University of Delaware
\item Samuel Fletcher (UC Irvine): University of Minnesota
\item Nicholas Teh (Cambridge): University of Notre Dame
\item Joseph Rachiele*: Chapman University
\item Joshua Hershey: The King's College NY
\item Jo E. Wolff (Stanford): University of Edinburgh
\item Jim Weatherall (UC Irvine): UC Irvine
\item Caleb Cohoe: Metropolitan State University, Denver
\item Jada Strabbing: Fordham University, Wayne State University
\item Giovanni Valente (Maryland): University of Pittsburgh, Politecnico Milano
\item Yuichiro Kitajima (Hokkaido): Nihon University
\item David J.\ Baker*: University of Michigan
\item Lara Buchak: UC Berkeley, Princeton
\item Tracy Lupher (Texas): James Madison University
\item Antony Eagle: Oxford University, University of Adelaide
\end{itemize}

Students in Princeton's PhD program earn ``units of credit'' by writing papers
under the supervision of a member of the faculty. I have advised unit work in
general philosophy of science, philosophy of physics, philosophy of religion,
and philosophical logic.

\rubric{Graduate seminars}

\begin{itemize}
\item Spacetime and objectivity
\item Kierkegaard's \textit{Concluding Unscientific Postscript}
\item Philosophy of space and time (with D.\ Hogan)
\item Logical philosophy of science
\item Foundations of mathematics: set theory vs.\ category theory (with J.\ Burgess)
\item Categorical logic and topos theory
\item Quantum information theory and the foundations of quantum mechanics
  (with B.\ van Fraassen)
\item Metaphysics of physics (with J.\ Butterfield)
\item Foundations of quantum field theory
\item Scientific realism and antirealism
\item From physics to metaphysics (with A.\ Elga)
\end{itemize}

\rubric{Undergraduate courses and seminars}

\begin{itemize}
\item Religion and Scientific Objectivity
\item Kierkegaard
\item Philosophy of Science
\item History of Analytic Philosophy 1900--1950
\item Philosophy of Mathematics
\item Category Theory (math--phil)
\item Advanced Logic (math--phil)
\item Intermediate Logic
\item Philosophical Logic
\item Philosophy of Physics
\item Philosophy of Religion
\item Introductory Logic
\item History of Modern Philosophy
\item The Modern Breakthrough in Scandinavia
\end{itemize}

\rubric{Senior theses}

Since 2001, between one and five theses per year, for Philosophy,
Mathematics, Physics, and ORFE. Besides my areas of expertise, I have advised
theses in pragmatism, metaphysics, philosophy of mind, religion, biology, and
medicine.


\subsection*{Academic service}

\rubric{University and departmental administration}

\begin{itemize}
\item Princeton Committee on Undergraduate Admission and Financial Aid,
  2024--27
\item Princeton Policy Committee on Athletics and Campus Recreation, 2022--25
\item Princeton University Research Board, 2007--08, 2009--11, 2016--17
\item Department of Philosophy:
  \begin{itemize}
  \item Director of Graduate Studies, 2007--08
  \item appointments committee, 2018--20
  \item director of PhD placement, 2009--12, 2016--17, 2022--
  \item director of colloquium committee, 2010--11, 2016--17
  \item graduate admissions committee, 2001--02, 2003--06, 2007--08,
    2009--11, 2013--14
  \item curriculum committee, 2001--03, 2004--06, 2007--08
  \item computer committee, 2004--06, 2014--15
  \item library committee, 2009--10
  \item graduate language examiner (German), 2001--
  \end{itemize}
\end{itemize}

\rubric{Editorial work}

\begin{itemize}
\item Associate editor, \emph{Philosophy of Physics}, 2022--
\item Associate editor, \emph{Review of Symbolic Logic}, 2020--
\item Associate editor, \emph{Science, Religion and Culture}, 2017--
\item Area editor for philosophy of physical science,
  \href{http://www.philpapers.org}{PhilPapers.org}, 2011--
\end{itemize}

\rubric{Conference organization and program committees}

\begin{itemize}
\item Co-organizer, Oxford--Princeton partnership in philosophy of physics,
  2001--07
\item Organizer, Conference on understanding quantum mechanics through
  mathematical innovation, Princeton, Oct 2007
\item Co-organizer, Symposium on categorical perspectives in the philosophy of
  physics, Philosophy of Science Association Meeting, 2010
\item Co-organizer, Irvine--Pittsburgh--Princeton conference series on the
  foundations of physics, 2013--15
\item Organizer, Conference on equivalent theories in science and metaphysics,
  Princeton, Mar 2015
\item Co-organizer, Foundations of Quantum Theory Conference, Nagoya, Japan,
  Mar 2015
\item Co-organizer, Workshop on Georg Brandes and the Modern Breakthrough,
  May 2025
\item Co-organizer, Workshop on Kierkegaard, Objectivity, and Subjectivity,
  Jul 2025
\item Program committee:
  \begin{itemize}
  \item Philosophy of Science Association Meeting, 2006
  \item The semantics of scientific theories, LMU Munich, 2016
  \item European Philosophy of Science Association Meeting, 2023
  \item Foundations of Physics, 2026
  \end{itemize}
\end{itemize}

\rubric{External evaluation and advisory work}

\begin{itemize}
\item Scientific Advisory Board, ERC Advanced Grant project ``Team semantics
  and dependence'' (PI: Jouko V{\"a}{\"a}n{\"a}nen), 2021--24
\item External reviewer for faculty hiring and promotion: Boston University,
  Chicago, Columbia, CUNY, CCNY, Drexel, Georgetown, London School of
  Economics, National Technical University of Athens, Notre Dame, NYU, Oxford,
  Radboud University of Nijmegen, Rutgers, St.\ Olaf College, Stanford, UC
  Irvine, University of Chicago, University of Edinburgh, University of
  Michigan, Pittsburgh, University of Southern California, University of
  Vienna, University of Washington, University of Waterloo, Wheaton College,
  Yale
\item External dissertation examiner: University of Geneva; University of
  Texas; University of Utrecht (foundations of science); Radboud University of
  Nijmegen (mathematics); Carnegie Mellon University; LMU Munich; University
  of Aberdeen (theology); Roskilde Universitet; University of Salzburg
\item Academic Advisory Board, National Speech and Debate Association, 2015--16
\item Invited participant, New York Institute of Philosophy, cosmology and
  religion discussion group, 2008
\item Invited participant, Princeton Council of the Humanities, philosophy of
  time group, 2003--04
\end{itemize}

\rubric{Refereeing}

\begin{itemize}
\item Journals and conferences: \journal{Analysis}, \journal{Annals of
    Physics}, \journal{Annals of Pure and Applied Logic},
  \journal{Australasian Journal of Philosophy}, \journal{British Journal for
    the Philosophy of Science}, \journal{Bulletin of Symbolic Logic},
  \journal{Canadian Journal of Philosophy}, \journal{Christian Scholars
    Review}, \journal{Communications in Mathematical Physics},
  \journal{Erkenntnis}, \journal{European Journal for Philosophy of Science},
  \journal{Foundations of Physics}, \journal{International Journal for
    Philosophy of Religion}, \journal{International Studies in the Philosophy
    of Science}, \journal{Journal for General Philosophy of Science},
  \journal{Journal of Logic \& Analysis}, \journal{Journal of Mathematical
    Physics}, \journal{Journal of Philosophy}, \journal{Journal of Physics A},
  \journal{Logic and Logical Philosophy}, \journal{Notre Dame Journal of
    Formal Logic}, \journal{Nous}, \journal{Philosophia Mathematica},
  \journal{Philosopher's Imprint}, \journal{Philosophical Studies},
  \journal{Philosophical Transactions of the Royal Society},
  \journal{Philosophy and Phenomenological Research}, \journal{Philosophy
    Compass}, \journal{Philosophy of Science}, \journal{Physical Review~A},
  \journal{Physical Review Letters}, \journal{Physics Letters~A},
  \journal{Quantum Studies: Mathematics and Foundations}, \journal{Ratio},
  \journal{Religious Studies}, \journal{Reviews of Mathematical Physics},
  \journal{Science and Christian Belief}, \journal{Studia Logica},
  \journal{Studies in History and Philosophy of Modern Physics},
  \journal{Studies in History and Philosophy of Science}, \journal{Theoria};
  Quantum Physics and Logic (QPL 2011).
\item Funding agencies: National Science Foundation (USA); Engineering and
  Physical Sciences Research Council (UK); Israel Science Foundation;
  Netherlands Organisation for Scientific Research (NWO); Royal Netherlands
  Academy of Arts and Sciences (KNAW); Social Sciences and Humanities Research
  Council of Canada; Agence Nationale de la Recherche (France); Austrian
  Science Fund (FWF); German-Israeli Foundation for Scientific Research and
  Development; European Research Council; National Research, Development and
  Innovation Office of Hungary; Foundation for Polish Science; The Leverhulme
  Trust; John Templeton Foundation; Templeton World Charity Foundation.
\item Publishers: Elsevier; Hackett; Oneworld; Oxford University Press;
  Princeton University Press; Springer; Wiley-Blackwell; Yale University
  Press.
\end{itemize}


\subsection*{Languages}

\begin{itemize}
\item English, native speaker
\item Danish, C1 (Studieprøven)
\item German, C1 (Zentrale Mittelstufenprüfung)
\item Dutch, B2
\item French, reading
\end{itemize}


\vspace{2em} \hfill \textit{Last updated: \today}

\end{document}

%%% Local Variables:
%%% mode: latex
%%% TeX-master: t
%%% End:
